% Vietnamese translation for OI Public Library
% Source-PDF: 英文资料 ENGLISH MATERIALS/Concrete Mathematics - 2nd Edition.pdf
\documentclass[11pt]{article}
\usepackage{oipl-vn}

\title{Toán học cụ thể}
\author{Ronald L. Graham, Donald E. Knuth, Oren Patashnik}
\date{Addison-Wesley, ấn bản thứ hai, 1994}

\begin{document}
\maketitle

\origtitle{Concrete Mathematics: A Foundation for Computer Science}
\sourcepath{English Materials / Concrete Mathematics - 2nd Edition.pdf}

\renewcommand{\abstractname}{Tóm tắt}
\begin{abstract}
Đây là ghi chú dịch tiếng Việt cho \emph{Concrete Mathematics}, một sách nền
tảng về các kỹ thuật toán học rời rạc dùng trong khoa học máy tính. Bản PDF có
lớp văn bản đọc được nhưng một số ký tự trong OCR và phông toán học bị sai.
Ghi chú này chuyển ngữ thông tin đầu sách, cấu trúc chương và vai trò của từng
phần đối với luyện thuật toán.
\end{abstract}

\section{Thông tin sách}

\begin{center}
\begin{tabular}{ll}
\textbf{Tác giả} & Ronald L. Graham; Donald E. Knuth; Oren Patashnik \\
\textbf{Nhà xuất bản} & Addison-Wesley \\
\textbf{Ấn bản} & Thứ hai \\
\textbf{Năm} & 1994 \\
\textbf{ISBN} & 0-201-55802-5
\end{tabular}
\end{center}

Sách được viết từ khóa học cùng tên tại Stanford. Mục tiêu là rèn khả năng
biến đổi công thức, giải truy hồi, tính tổng hữu hạn, dùng hàm sinh và ước
lượng tiệm cận. Đây là tài liệu nền tảng rất sát với các kỹ thuật thường gặp
trong phân tích thuật toán và tổ hợp tính toán.

\section{Mục lục chương đã dịch}

\begin{description}
  \item[Chương 1. Các bài toán truy hồi]
  Dẫn nhập qua các ví dụ cổ điển như số miền do đường thẳng tạo ra và bài toán
  Josephus; giới thiệu cách đoán, khai triển và chứng minh nghiệm truy hồi.

  \item[Chương 2. Tổng]
  Các phép biến đổi tổng hữu hạn, đổi thứ tự tổng, tổng kép, tổng dạng hữu hạn
  và các kỹ thuật chuẩn để biến tổng phức tạp thành biểu thức đóng.

  \item[Chương 3. Hàm nguyên]
  Hàm sàn, hàm trần, phần dư, phần phân số và các đồng nhất thức liên quan;
  đây là phần rất gần với bài toán chia đoạn, chia khối và đếm theo thương.

  \item[Chương 4. Lý thuyết số]
  Chia hết, số nguyên tố, ước chung, thuật toán Euclid, đồng dư, hàm số học và
  các công cụ số học rời rạc thường dùng trong thuật toán.

  \item[Chương 5. Hệ số nhị thức]
  Các đồng nhất thức nhị thức, biến đổi tổ hợp, tổng chứa hệ số nhị thức và
  phương pháp dùng chúng trong đếm.

  \item[Chương 6. Các số đặc biệt]
  Số Stirling, số Euler, số harmonic, Bernoulli và các họ số xuất hiện tự nhiên
  trong đếm, hoán vị, phân hoạch và khai triển.

  \item[Chương 7. Hàm sinh]
  Hàm sinh thường, hàm sinh mũ, thao tác đại số trên chuỗi hình thức và cách
  dùng hàm sinh để giải truy hồi hoặc bài toán đếm.

  \item[Chương 8. Xác suất rời rạc]
  Kỳ vọng, biến ngẫu nhiên rời rạc, phân phối cơ bản và ứng dụng xác suất trong
  đếm cũng như phân tích thuật toán.

  \item[Chương 9. Tiệm cận]
  Ký hiệu tiệm cận, ước lượng tổng, xấp xỉ hàm, phương pháp phân tích hành vi
  khi \(n\) lớn và các công cụ cần cho đánh giá độ phức tạp.
\end{description}

\section{Ghi chú nội dung}

Trong bối cảnh OI và lập trình thi đấu, các phần hữu ích nhất thường là:
\begin{itemize}
  \item Chương 1, 2 và 7 cho truy hồi, tổng hữu hạn và hàm sinh.
  \item Chương 3 và 4 cho các bài toán số học, chia theo thương và đồng dư.
  \item Chương 5 và 6 cho tổ hợp nâng cao và các dãy số đặc biệt.
  \item Chương 8 và 9 cho phân tích thuật toán ngẫu nhiên và ước lượng tiệm cận.
\end{itemize}

\section{Ghi chú về trích xuất}

Tệp PDF có 670 trang. Văn bản tiếng Anh trích xuất được, nhưng bản PDF sinh từ
DVI cũ nên một số dấu ngoặc, dấu nối, ký hiệu toán và chữ trong lề bị OCR sai
hoặc mất mã hóa. Bản dịch đầy đủ từng chương cần đối chiếu trực tiếp PDF gốc
khi chuyển công thức, bảng và ghi chú lề sang LaTeX.

\end{document}
